\begin{enumerate}[label=\textbf{1.1.\arabic*.},ref=1.1.\arabic*]
\setlength{\itemsep}{0.35em}
\item \label{formula-1-1-1} \textbf{Matrices.} The transpose $A^{\mathsf T}$ is obtained by interchanging rows and columns. A square matrix is invertible exactly when $\det A\ne0$, and then $A^{-1}=\operatorname{adj}(A)/\det A$. A system $AX=B$ has solution $X=A^{-1}B$ when $A^{-1}$ exists.
\item \label{formula-1-1-2} \textbf{Coordinate geometry.} For points $A(x_1,y_1)$ and $B(x_2,y_2)$,
\begin{equation*}\tag{1.1.2.1}
AB=\sqrt{(x_2-x_1)^2+(y_2-y_1)^2},\qquad
m_{AB}=\frac{y_2-y_1}{x_2-x_1}.
\end{equation*}
Three points are collinear when the slopes of two connecting segments are equal.
\item \label{formula-1-1-3} \textbf{Limits and derivatives.}
\begin{equation*}\tag{1.1.3.1}
\lim_{x\to0}\frac{\sin x}{x}=1,\qquad
\frac{d}{dx}e^{g(x)}=e^{g(x)}g'(x),\qquad
\frac{d}{dx}x^x=x^x(1+\ln x)\quad(x>0).
\end{equation*}
\item \label{formula-1-1-4} \textbf{Inverse-trigonometric integral.}
\begin{equation*}\tag{1.1.4.1}
\int\frac{dx}{\sqrt{a^2-b^2x^2}}=\frac1b\sin^{-1}\!\left(\frac{bx}{a}\right)+C.
\end{equation*}
For a bounded region under $y=f(x)$, its area is $\int_a^b f(x)\,dx$ when $f(x)\ge0$ on $[a,b]$.
\item \label{formula-1-1-5} \textbf{Linear programming.} The maximum or minimum of a linear objective over a polygonal feasible region occurs at a feasible vertex; if the objective line is parallel to an optimal boundary edge, every point on that edge is optimal.
\item \label{formula-1-1-6} \textbf{Vectors.}
\begin{equation*}\tag{1.1.6.1}
\vec a\cdot\vec b=|\vec a||\vec b|\cos\theta,\qquad
\operatorname{Area}(\triangle ABC)=\frac12|\overrightarrow{AB}\times\overrightarrow{AC}|.
\end{equation*}
\item \label{formula-1-1-7} \textbf{Lines in 3D.} A line through $\vec a$ in direction $\vec d$ is $\vec r=\vec a+t\vec d$. The shortest distance between skew lines is
\begin{equation*}\tag{1.1.7.1}
d=\frac{|(\vec a_2-\vec a_1)\cdot(\vec d_1\times\vec d_2)|}{|\vec d_1\times\vec d_2|}.
\end{equation*}
\item \label{formula-1-1-8} \textbf{Probability.} $P(E^c)=1-P(E)$; for independent trials, $P(\text{at least one success})=1-(1-p)^n$. The total-probability and Bayes formulas are
\begin{equation*}\tag{1.1.8.1}
P(B)=\sum_iP(B\mid A_i)P(A_i),\qquad
P(A_i\mid B)=\frac{P(B\mid A_i)P(A_i)}{P(B)}.
\end{equation*}
\item \label{formula-1-1-9} \textbf{First-order linear differential equations.} For $y'+P(x)y=Q(x)$, the integrating factor is $\mu(x)=e^{\int P(x)\,dx}$ and $(\mu y)'=\mu Q$.
\item \label{formula-1-1-10} \textbf{Chain rule.} If $y=F(u(x))$, then $dy/dx=(dF/du)(du/dx)$; for differentiation of $F$ with respect to $G$, use $(dF/dx)/(dG/dx)$ where $dG/dx\ne0$.
\end{enumerate}
